This project stems from the initiative of the interest group Quantum@UC. This group is a multidisciplinary community made up of individuals interested in quantum technologies, with the shared goal of developing and exploring new possibilities in science and engineering in the fields of quantum computing, communication, and sensing at the University of Coimbra.

The Quantum@UC group emerged from the academic community’s interest in quantum technologies and from the need to further develop this field at the University of Coimbra. In the Department of Physics and the Department of Electrical and Computer Engineering, optional courses in quantum technologies and quantum computing were created by Professor Helena Alberto and Professor Jorge Lobo, respectively. In collaboration with the Department of Informatics Engineering, the group was formed, with Professor Edmundo Monteiro appointed by the University’s Rectorate to coordinate it. Since then, the group has organized several initiatives, including the creation of a non-degree course in quantum technologies for the 2024–2025 academic year \cite{QuantumUC_about}.

The project itself is being developed at QuantumLab@UC, a new research laboratory at the University of Coimbra focused on quantum technologies, funded by the Recovery and Resilience Plan (PRR) and the European Union. The laboratory has two rooms, one in the Department of Electrical and Computer Engineering and the other in the Department of Physics, equipped with two NMR-based quantum computers, one with two qubits and the other with three qubits, as well as experimental Quantum Key Distribution (QKD) equipment for demonstrating the principles of quantum cryptography.

For the development of this project, I am being supervised by Professor Sagar Pratapsi and co-supervised by Professor Rui Vilão, both members of the Quantum@UC group.

The project arises as a continuation of the study carried out by Professor Sagar Pratapsi et al. on the characterization of Bargmann invariants \cite{Pratapsi2025elementary}, with the measurement of this quantity being the proposed procedure for testing the practical limitations of the NMR-based quantum computer.